Event detection turns long time series into discrete decisions: alarms, boundaries, actions, or episodes that are scored by their timing.
A common reduction trains samplewise segmentation masks and derives event times only after thresholding or transition detection, creating a mismatch when evaluation ranks temporal detections rather than labels at every sample.
We introduce Boundary Density Likelihood (BDL), a likelihood objective that trains sequence models on the events later matched by event average precision.
BDL maps annotations to expected events on the model output timeline and fits nonnegative sequence outputs with a Poisson log score.
Unit-sum smoothing represents timing uncertainty, while summing into output bins changes temporal resolution without changing the number of annotated events.
Peaks in the resulting score sequences form ranked detections for standard tolerance-based matching.
Across two benchmarks, BDL improves event AP under matched detectors: on sleep, mAP improves from 0.575 to 0.680 over cross-entropy segmentation and one-minute AP improves from 0.051 to 0.219; on CHB-MIT, one-second-stride seizure mAP improves from 0.287 to 0.388.
Controlled ablations show that the gains are not explained by class imbalance, smoothing, or a single backbone.
Additional point-event, stride, streaming-context, calibration, and state-reconstruction studies identify where the objective helps and where detector design still matters.
